\section{Elementary convex facts used in the proof}\label{app:convex-facts}
This appendix proves the finite-dimensional convex facts needed for the
pressure limits and the persistence argument. These are special cases of the
convex convergence, conjugacy, and reachable-gradient results in
\cite[Corollary~7.18 and Theorems~11.1 and 9.61]{RockafellarWets1998}. The proofs below
make their exact hypotheses explicit. Lebesgue measure and its Fubini theorem
are used only to locate ordinary differentiability points of a planar convex
function.

\begin{lemma}[Local convex convergence]\label{lem:elementary-convex-convergence}
Let $O\subset\R^m$ be open and convex. A finite convex function $F$ on $O$
is locally Lipschitz and has a nonempty compact subdifferential at every
point of $O$. Its subgradients are bounded on compact subsets. If finite
convex functions $F_n$ converge pointwise on $O$ to a finite function $F$,
then $F$ is convex and the convergence is locally uniform. Moreover, if
$x_n\to x\in O$, $F_n\to F$ locally uniformly, and
$v_n\in\partial F_n(x_n)$, then $(v_n)$ is bounded and every limit belongs
to $\partial F(x)$.
\end{lemma}

\begin{proof}
Choose a cube with center $c$ whose concentric enlargement lies in $O$.
Convexity bounds a function above on the larger cube by its values at the
vertices. The inequality
$F(y)\ge 2F(c)-F(2c-y)$ then bounds it below on a smaller concentric cube.
On a still smaller cube, extend any chord by a fixed positive distance in
both directions while remaining in the bounded region. Monotonicity of
one-dimensional convex secant slopes bounds the difference quotient on the
original chord by its surrounding oscillation divided by that distance.
This proves local Lipschitz continuity. In the sequence assertion, the
finitely many vertex values and the center value are uniformly bounded by
pointwise convergence, so the same argument gives a common local Lipschitz
constant. A finite net of points on each compact set upgrades pointwise to
uniform convergence. Passing convexity inequalities to limits proves
convexity of $F$.

Here is also a direct supporting-plane construction. Separate a point
strictly below the epigraph from that convex epigraph, let the point tend to
$(x,F(x))$, and normalize the separating normals. Separation in finite
dimensions follows by projecting an exterior point onto a closed convex set:
if $z$ is its nearest point, expansion of the squared distance along
$z+\theta(y-z)$ gives the separating inequality. One can restrict the
epigraph to a closed neighborhood of $x$ before applying this construction.
A limiting normal has a nonzero vertical component: a horizontal supporting
normal would separate $x$ from an open neighborhood of $x$, which is
impossible. Division by that component produces a local affine minorant
through $(x,F(x))$. Convexity extends it to all of $O$ by considering the
initial portion of the segment from $x$ to any other point. Thus
$\partial F(x)$ is nonempty.

If $v\in\partial F(x)$ and $x\pm re_j\in O$, the supporting inequalities
bound $|v\cdot e_j|$ by the oscillation of $F$ divided by $r$.
Uniform oscillation on a surrounding compact set bounds all subgradients
at points in a smaller compact set. Subdifferentials are closed by their
defining inequalities, hence compact. The same bounds apply to $F_n$ under
local uniform convergence. Passing
$F_n(y)\ge F_n(x_n)+v_n\cdot(y-x_n)$ to a convergent subsequence proves the
last assertion.
\end{proof}

\begin{lemma}[One-dimensional conjugacy]\label{lem:elementary-conjugacy}
Let $u:\R\to\R\cup\{+\infty\}$ be proper, lower semicontinuous, and convex,
with a domain having nonempty interior. Then
\[
 u(d)=\sup_{h\in\R}\{hd-u^*(h)\},\qquad
 u^*(h)=\sup_{d\in\R}\{hd-u(d)\}.
\]
At a finite point $d$, equality $u(d)+u^*(h)=hd$ holds precisely when
$u(d')\ge u(d)+h(d'-d)$ for every $d'$. In particular, if $K\subset\R^2$
is nonempty compact convex with a density projection of positive length,
then for every $x\in\R$,
\[
 \max\{xs:(s,d)\in K\}
 =\inf_{z\in\R}\{h_K(x,z)-zd\},
\]
where the maximum is $-\infty$ outside that projection.
\end{lemma}

\begin{proof}
At every interior point of the domain, a slope between the left and right
convex derivatives gives an affine minorant. Hence $u$ has at least one
affine minorant $\ell_0$. To obtain enough such minorants, fix a point
$(d,a)$ strictly below the closed epigraph of $u$. Strict separation, by
the projection argument in the preceding proof, yields an affine inequality
\[
 A d'+B y\ge c>A d+B a\qquad(y\ge u(d')).
\]
Necessarily $B\ge0$. If $B>0$, division gives an affine minorant of $u$
whose value at $d$ exceeds $a$. If $B=0$, add a sufficiently small positive
multiple of the valid inequality $y\ge\ell_0(d')$; the resulting
inequality has positive vertical coefficient and still strictly separates
$(d,a)$. This also gives an affine minorant exceeding $a$ at $d$.
Since every point below the epigraph is so separated, $u$ is the supremum
of its affine minorants, including at boundary and exterior points of its
domain. Among affine minorants of slope $h$, the largest is
$d\mapsto hd-u^*(h)$, which proves biconjugacy. Rearranging the definition
of $u^*$ proves the equality criterion.

For the last assertion, put
$k_x(d)=\max\{xs:(s,d)\in K\}$. Compactness and convexity of $K$ make
$k_x$ upper semicontinuous and concave on its compact projection interval.
Thus $-k_x$, extended by $+\infty$, satisfies the preceding hypotheses.
The identity $h_K(x,z)=\sup_d(k_x(d)+zd)$ and biconjugacy give the formula.
\end{proof}

\begin{lemma}[Directional and reachable gradients]\label{lem:reachable-local}
Let $F$ be finite and convex on an open convex subset $O\subset\R^2$.
For $x\in O$, let $G_F(x)$ be the limits of gradients at differentiability
points tending to $x$. Then $G_F(x)$ is nonempty compact and
\[
 \partial F(x)=\conv G_F(x).
\]
For every $e\in\R^2$,
\[
 \lim_{a\downarrow0}\frac{F(x+ae)-F(x)}a
 =\max_{v\in\partial F(x)}v\cdot e.
\]
The rescaled functions in this last display converge uniformly for $e$ in
compact subsets of $\R^2$.
\end{lemma}

\begin{proof}
On every coordinate line, a one-dimensional convex function has equal left
and right derivatives except at countably many points: both derivatives
are monotone, and each strict gap contains a distinct rational number.
Fubini's theorem therefore implies that both coordinate partial derivatives
of $F$ exist almost everywhere. At such a point, every subgradient has these
two coordinates, so the subdifferential is a singleton, say $\{v\}$.
This implies differentiability. Indeed, for small $u$ choose
$v_u\in\partial F(x+u)$. Local boundedness and closedness of subgradients
show $v_u\to v$, and the two supporting inequalities give
\[
 v\cdot u\le F(x+u)-F(x)\le v_u\cdot u.
\]
Their difference is $o(|u|)$. Thus differentiability points have full
measure, in particular they are dense.

Local boundedness and diagonal subsequences show that $G_F(x)$ is nonempty
and compact. Each of its elements is in $\partial F(x)$ by the limiting
support inequality. Conversely, take $p\in\partial F(x)$ and a direction
$e$. Choose differentiability points
$y_j=x+a_je+o(a_j)$, with $a_j\downarrow0$. This is possible by density,
for example with error less than $a_j^2$. Adding the supporting
inequalities at $x$ and $y_j$ gives
\[
 (\nabla F(y_j)-p)\cdot(y_j-x)\ge0.
\]
After passing to a bounded-gradient subsequence and dividing by $a_j$,
one obtains $v\in G_F(x)$ with $v\cdot e\ge p\cdot e$.
Thus $G_F(x)$ and $\partial F(x)$ have the same support function.
In the plane the convex hull of a compact set is compact: any convex
combination can be reduced to three terms by eliminating an affine
dependence among four or more points. Separation therefore proves
$\partial F(x)=\conv G_F(x)$.

Finally, the convex secant slope in the last assertion has a limit as
$a\downarrow0$. Supporting inequalities at $x$ give the lower bound
$v\cdot e$ for every $v\in\partial F(x)$. Choose instead
$v_a\in\partial F(x+ae)$; the reverse supporting inequality bounds the
same secant slope above by $v_a\cdot e$. A bounded subsequence of $v_a$
converges to an element of $\partial F(x)$, proving the claimed equality.
The local Lipschitz constant for $F$ gives a common Lipschitz constant
in $e$ for all sufficiently small rescalings on any fixed compact set.
A finite-net argument then proves uniform convergence there.
\end{proof}
